% ---------------------------------------------------------------- remerciements
\chapter*{Remerciements}
\markboth{Remerciements}{}

Je tiens à remercier Monsieur \textbf{Hamza Abderrahmen}, directeur général de
\textbf{CyberMech Systems}, de m'avoir accueillie au sein du département R\&D / Robotique
pour ce stage d'été.

J'exprime ma gratitude à \acompleter{Nom de l'encadrant(e)}, mon encadrant(e) au sein de
CyberMech Systems, pour son accompagnement tout au long du stage, sa disponibilité et la
pertinence de ses retours.

Mes remerciements s'adressent aussi à l'ensemble de l'équipe du projet AMR-X, avec qui
j'ai eu le plaisir de travailler, ainsi qu'au corps professoral et administratif de l'INSAT
pour la formation reçue.

% ---------------------------------------------------------------- tables
\cleardoublepage
\renewcommand{\contentsname}{Table des matières}
\tableofcontents
\cleardoublepage
\renewcommand{\listfigurename}{Table des figures}
\listoffigures
\cleardoublepage
\renewcommand{\listtablename}{Liste des tableaux}
\listoftables

% ---------------------------------------------------------------- acronymes
\cleardoublepage
\chapter*{Liste des acronymes}
\addcontentsline{toc}{chapter}{Liste des acronymes}
\markboth{Liste des acronymes}{}
\begin{tabular}{@{}p{3cm}l@{}}
  \textbf{AMCL}   & Adaptive Monte Carlo Localization \\[0.3em]
  \textbf{AMR}    & Autonomous Mobile Robot \\[0.3em]
  \textbf{API}    & Application Programming Interface \\[0.3em]
  \textbf{JSON}   & JavaScript Object Notation \\[0.3em]
  \textbf{MPPI}   & Model Predictive Path Integral \\[0.3em]
  \textbf{Nav2}   & Navigation 2 (pile de navigation ROS~2) \\[0.3em]
  \textbf{QoS}    & Quality of Service \\[0.3em]
  \textbf{REST}   & Representational State Transfer \\[0.3em]
  \textbf{ROS 2}  & Robot Operating System 2 \\[0.3em]
  \textbf{RViz}   & ROS Visualization \\[0.3em]
  \textbf{SDF}    & Simulation Description Format \\[0.3em]
  \textbf{SLAM}   & Simultaneous Localization and Mapping \\[0.3em]
  \textbf{SPA}    & Single Page Application \\[0.3em]
  \textbf{SVG}    & Scalable Vector Graphics \\[0.3em]
  \textbf{TF}     & Transform (arbre des repères ROS) \\[0.3em]
  \textbf{UI}     & User Interface \\[0.3em]
\end{tabular}

% ---------------------------------------------------------------- glossaire
\cleardoublepage
\chapter*{Glossaire}
\addcontentsline{toc}{chapter}{Glossaire}
\markboth{Glossaire}{}
\begin{description}[style=nextline,leftmargin=4.8cm,labelwidth=4.5cm,font=\bfseries]
  \item[Acquittement] Réponse d'un service ROS indiquant que la demande a été reçue et
    acceptée ou refusée. Un acquittement positif ne signifie pas que la destination est
    atteinte.
  \item[Costmap] Grille d'occupation utilisée par Nav2, dont chaque cellule contient un
    coût de 0 (libre) à 100 (occupé) ou $-1$ (inconnu).
  \item[Dashboard] Application web existante du projet AMR-X, composée d'un frontend React
    et d'un backend FastAPI, servant à la supervision de la flotte.
  \item[Donnée périmée] État de navigation reçu depuis plus de 2,5~s ; le dashboard
    le considère comme non fiable et bloque toute nouvelle commande.
  \item[Monde SDF] Description de l'environnement simulé chargée par Gazebo ; le dashboard
    en extrait un plan 2D pour l'affichage.
  \item[Passerelle] Nœud ROS~2 \code{dashboard\_navigation} développé pendant le stage,
    intermédiaire unique entre le navigateur et Nav2.
  \item[Phase] État courant d'une demande de navigation :\\ \code{idle}, \code{sending},
    \code{navigating}, \code{canceling}, \code{succeeded},
    \code{canceled}, \code{failed} ou \code{rejected}.
  \item[Repère \code{map}] Repère global de Nav2, défini par la carte chargée. Il ne coïncide
    pas nécessairement avec le repère du monde Gazebo.
  \item[rosbridge] Serveur exposant les topics et services ROS~2 au navigateur à travers
    une connexion WebSocket (port 9090).
\end{description}
